\section{Causas del Problema}
En el párrafo dos del capítulo anterior, mencionamos una gran cantidad de problemas dentro de la defensoría, mismos que eran causantes de un problema que afecta el proceso de una queja, el cual era el tiempo, sin embargo, existen circunstancias que causan cada uno de estos problemas, por ejemplo:

\begin{itemize}
    \item \textbf{Problemas de organización:} Encontramos que hay una gran rotación de personal, impidiendo la adaptación al trabajo, pues al no existir un plan de trabajo específico, provoca que la adaptabilidad al programa se tome una gran cantidad de tiempo, mismo que podría ser usado para la resolución de quejas.
    
    \item \textbf{Deficiencias en el manejo de expedientes:} Encontramos trabas en la comunicación, pues todo se rige a través de Excels compartidos donde no existe una nomenclatura específica para el registro, lo que trae consigo una pérdida de información al querer buscar expedientes en específico, así como la duplicidad de folios o números de expedientes.
    
    \item \textbf{Falta de personal:} Está ligada directamente al IPN, situaciones que desconocemos, sin embargo, es claro que a mayor personal, mayor respuesta.
    
    \item \textbf{Problemas del quejoso:} Son por la falta de transparencia con sus quejas, pues al no tener una respuesta pronta y además desconocer la situación de su queja, provoca una duplicidad en las mismas.
\end{itemize}

Entonces, la falta de tiempo tiene una gran cantidad de problemas raíz, y el unificar su proceso a un sistema, aún con sus carencias, mejoraría el tiempo de respuesta, atacando directamente el problema general.